\documentclass{article}
\usepackage[margin=0.5in, left=0.5in, right=0.5in]{geometry}
\usepackage{amsmath}

\title{Assignment 1 Report \\ Digital Signal Processing \\ Fourier Transform}
\author{James Middleton (2785209M) \& Ewan Stewart (2664680S)}
\date{\today}

\begin{document}
\maketitle

\section*{Introduction}
We have developed software to transform, filter, and modulate audio signals utilising Fourier transforms and the frequency domain.

\section{Audio Plot}
Part 1 contains the logic for obtaining the frequency domain representation of our audio signal, involving reading the WAV file, normalizing the signal, performing a fast Fourier transform, and plotting the graphs in the time and frequency domains.

\subsection{Read\_WAV\_file}
Reads and normalizes the data to a range of [-1:1] (32-bit float). If a stereo input is given, only the first channel is stored. Automatically handles 8-bit unsigned, 16-bit signed, and 32-bit signed integer PCM.
\begin{description}
    \item[Inputs:]
    \item[Outputs:] 
\end{description}

\subsection{Plot\_time\_spectrum}
Plots a time domain graph representing the audio signal (amplitude against time).
\begin{description}
    \item[Inputs:]
    \item[Outputs:] 
\end{description}

\subsection{calculate\_fft}
Takes a fast Fourier transform of the audio data, producing a full complex spectrum.
\begin{description}
    \item[Inputs:]
    \item[Outputs:] 
\end{description}

\subsection{get\_nyquist\_spectrum}
Produces a Nyquist-adjusted frequency spectrum for plotting purposes.
\begin{description}
    \item[Inputs:]
    \item[Outputs:] 
\end{description}

\subsection{plot\_frequency\_spectrum}
Plots a frequency domain \.svg plot in the specified output folder, with optional logarithmic scaling.
\begin{description}
    \item[Inputs:]
    \item[Outputs:] 
\end{description}


\section{Audio Enhancer}
Frequency filtering and amplification.

\subsection{generate\_wav\_file}
Generates a WAV file, designed to process frequency-modulated audio data.
\begin{description}
    \item[Inputs:]
    \item[Outputs:] 
\end{description}

\subsection{calculate\_ifft}
Performs an inverse Fourier transform on the input complex spectrum to produce a reconstructed time-domain data array.
\begin{description}
    \item[Inputs:]
    \item[Outputs:] 
\end{description}

\subsection{raised\_cosine\_filter}
Uses a cosine smoothing method to modulate/remove specific frequency ranges without producing ringing from Gibb's phenomenon.
\begin{description}
    \item[Inputs:]
    \item[Outputs:] 
\end{description}


\section{Part 3 \- Aural Exciter}
Utilizes the functions from the previous two parts of the assignment to digitally recreate an aural exciter circuit as given in the specification. An inverse tan function is applied to a high-pass filtered portion of the input signal to introduce non-linear harmonic distortion. These newly generated high-frequency harmonics are then scaled and mixed back with the original dry audio, enhancing the signal's brightness and presence.

\end{document}